\section{Morfismos equidimensionales}
\label{section:IV.13}

\oldpage[IV-3]{187}

Esta sección está dedicada al estudio de la variación de la \emph{dimensión} de las fibras de un morfismo localmente de tipo finito $f:X\to Y$ (que ya intervino a propósito de la «fórmula de dimensión» en~\sref{IV.5.5} y~\sref{IV.5.6}).
Demostramos primero \sref{IV.13.1.3} que la función $x\mapsto\dim_x f^{-1}(f(x))$ es siempre \emph{semicontinua superiormente} en $X$ (teorema de semicontinuidad de Chevalley).
Estudiamos después más particularmente los morfismos, llamados «equidimensionales», para los cuales esta función es \emph{localmente constante}.
\oldpage[IV-3]{188}
Desafortunadamente, la noción de morfismo equidimensional no es estable por cambio de base; por ello, en muchas cuestiones es más conveniente trabajar con la noción de morfismo universalmente abierto, cuyo estudio constituye el objeto de \textsection\textsection~14 y~15.

\subsection{Teorema de semicontinuidad de Chevalley}
\label{subsection:IV.13.1}

\begin{lemma}[13.1.1]
\label{IV.13.1.1}
Sea $Y$ un preesquema irreducible localmente noetheriano, $X$ un preesquema irreducible y $f:X\to Y$ un morfismo dominante de tipo finito.
Sea $\xi$ (resp.\ $\eta$) el punto genérico de $X$ (resp.\ de $Y$), y sea $e=\dim(f^{-1}(\eta))$.
Entonces, para todo $x\in X$, todas las componentes irreducibles de $f^{-1}(f(x))$ son de dimensión $\geq e$.
\end{lemma}

\begin{proof}
La afirmación es inmediata cuando, para todo $y\in f(X)$, $\sh{O}_y$ es un anillo \emph{universalmente catenario}: en efecto, si $x$ es el punto genérico de una componente irreducible $Z$ de $f^{-1}(y)$, de \sref{IV.5.6.5}, combinado con \sref{IV.5.2.1}, se sigue que $e+\dim(\sh{O}_y)=\dim(Z)+\dim(\sh{O}_x)$; pero, por \sref[0]{0.16.3.9}, se tiene $\dim(\sh{O}_x)\leq\dim(\sh{O}_y)$, de donde la conclusión en este caso.

Reduzcamos ahora el caso general a este caso particular.
La cuestión es evidentemente local sobre $Y$ y, teniendo en cuenta \sref{IV.4.1.1.3}, también es local sobre $X$; podemos, por tanto, restringirnos al caso en que $Y=\Spec(A)$ y $X=\Spec(B)$ son afines e irreducibles, siendo $B$ una $A$-álgebra de tipo finito.
Además, por \sref{IV.1.5.4}, podemos suponer que $X$ e $Y$ son reducidos, de modo que $A$ y $B$ son dominios íntegros.
Como $f$ es dominante, $A$ es entonces un \emph{subanillo} de $B$.
Consideremos $A$ como límite inductivo de sus sub-$\bb{Z}$-álgebras de tipo finito; de \sref{IV.8.9.1} se sigue entonces que existe una tal subálgebra $A_0$ y una $A_0$-álgebra de tipo finito $B_0$ tales que $B=B_0\otimes_{A_0}A$.
Pongamos $Y_0=\Spec(A_0)$, $X_0=\Spec(B_0)$, y sea $f_0:X_0\to Y_0$ el morfismo correspondiente al homomorfismo $A_0\to B_0$, de modo que $f=f_0\times 1_Y$.
No es evidente \emph{a priori} que el preesquema $X_0$ sea íntegro, pero veremos que podemos reducirnos a este caso.
Sea $\eta_0$ el punto genérico de $Y_0$, de modo que, si $g$ denota el morfismo $Y\to Y_0$, se tiene $g(\eta)=\eta_0$; por transitividad de las fibras \sref[I]{I.3.6.4}, se tiene $f^{-1}(\eta)=f_0^{-1}(\eta_0)\otimes_{k(\eta_0)}k(\eta)$ y, como $f^{-1}(\eta)$ es irreducible por hipótesis, también lo es $f_0^{-1}(\eta_0)$ \sref{IV.4.4.1}.
Nuestra afirmación resultará entonces del lema siguiente:

\begin{lemma}[13.1.2]
\label{IV.13.1.2}
Sean $Y_0$, $Y$ dos preesquemas íntegros de puntos genéricos $\eta_0$, $\eta$, sea $g:Y\to Y_0$ un morfismo dominante, y sea $f_0:X_0\to Y_0$ un morfismo dominante tal que $f_0^{-1}(\eta_0)$ es irreducible.
Sea $X_0'$ la única componente irreducible de $X_0$ que corta a $f_0^{-1}(\eta_0)$ \EGAReadLink{0.2.1.8}{($\mathbf{0}_{\mathrm{I}}$,~2.1.8)}, y denotemos también por $X_0'$ el subpreesquema cerrado reducido de $X_0$ cuyo espacio subyacente es $X_0'$.
Supongamos que el preesquema $X=X_0\times_{Y_0}Y$ es íntegro; entonces $X$ es isomorfo a $X_0'\times_{Y_0}Y$.
\end{lemma}

En efecto, si $j_0:X_0'\hookrightarrow X_0$ es la inyección canónica, que es una inmersión cerrada, $j=j_0\times 1_Y:X_0'\times_{Y_0}Y\to X_0\times_{Y_0}Y=X$ es una inmersión cerrada.
Por otra parte, $X_0'$ contiene la fibra $f_0^{-1}(\eta_0)$, de modo que $X_0'\times_{Y_0}Y$ contiene a $f^{-1}(\eta)$; obsérvese también que $f^{-1}(\eta)$ no es vacío \sref[I]{I.3.4.7}, por lo que contiene el punto genérico de $X$; en consecuencia, la imagen de $j$ es necesariamente todo $X$.
Pero, como $X$ es íntegro, el único subpreesquema cerrado de $X$ cuyo espacio subyacente es $X$ es el propio $X$, de modo que $j$ es un isomorfismo.

Establecido este lema, podemos suponer, por tanto, que $X_0$ es íntegro; para todo $y_0\in Y_0$, $\sh{O}_{y_0}$ es una $\bb{Z}$-álgebra esencialmente de tipo finito, luego un anillo universalmente
\oldpage[IV-3]{189}
catenario \sref{IV.5.6.4}; en consecuencia, para todo $y_0\in f_0(X_0)$, todas las componentes irreducibles de $f_0^{-1}(y_0)$\footnote{La fuente impresa dice aquí $f_0^{-1}(y)$; el punto cuantificado y el argumento requieren $f_0^{-1}(y_0)$.} tienen dimensión $\geq e$, pues sabemos que $\dim(f_0^{-1}(\eta_0))=e$ por transitividad de las fibras \sref{IV.4.1.4}.
Para todo $y\in f(X)$, se tiene entonces $y_0=g(y)\in f_0(X_0)$ y, por transitividad de las fibras y \sref{IV.4.2.7}, esto completa la demostración.
\end{proof}

\begin{theorem}[13.1.3]
\label{IV.13.1.3}
\textnormal{(Chevalley).} ---
Sea $f:X\to Y$ un morfismo localmente de tipo finito.
Para todo entero $n$, el conjunto $\mathrm{F}_n(X)$ de los $x\in X$ tales que $\dim_x(f^{-1}(f(x)))\geq n$ es cerrado; en otras palabras, la función $x\mapsto\dim_x(f^{-1}(f(x)))$ es semicontinua superiormente en $X$.
\end{theorem}

\begin{proof}
I) Supongamos primero que $f$ es localmente de presentación finita.
La cuestión es evidentemente local sobre $X$ y sobre $Y$, y podemos suponer, por tanto, que $Y=\Spec(A)$, $X=\Spec(B)$ son afines, siendo $B$ una $A$-álgebra de presentación finita.
Sabemos entonces por \sref{IV.8.9.1} que existen un subanillo noetheriano $A_0$ de $A$ y una $A_0$-álgebra de tipo finito $B_0$ tales que, poniendo $Y_0=\Spec(A_0)$ y $X_0=\Spec(B_0)$, se tiene
\[
X=X_0\times_{Y_0}Y,\qquad f=f_0\times 1_Y,
\]
donde $f_0:X_0\to Y_0$ es el morfismo correspondiente al homomorfismo $A_0\to B_0$.
Sea $g:Y\to Y_0$ el morfismo correspondiente a la inyección canónica $A_0\to A$, sea $y$ un punto de $Y$, y pongamos $y_0=g(y)$; sabemos que $f^{-1}(y)=f_0^{-1}(y_0)\otimes_{k(y_0)}k(y)$, y de \sref{IV.4.2.7} se sigue que, si $p$ es la proyección canónica $f^{-1}(y)\to f_0^{-1}(y_0)$, las componentes irreducibles de $f^{-1}(y)$ son las componentes irreducibles de los espacios $p^{-1}(Z_0)$, donde $Z_0$ recorre el conjunto de las componentes irreducibles de $f_0^{-1}(y_0)$, y cada componente irreducible de $p^{-1}(Z_0)$ domina $Z_0$ y tiene dimensión $\dim(Z_0)$.
Teniendo en cuenta \sref[0]{0.14.1.5}, vemos, por tanto, que si $g':X\to X_0$ es la proyección canónica, $x$ un punto de $X$ y $x_0=g'(x)$, se tiene $\dim_x(f^{-1}(f(x)))=\dim_{x_0}(f_0^{-1}(f_0(x_0)))$, de donde $\mathrm{F}_n(X)={g'}^{-1}(\mathrm{F}_n(X_0))$, y se reduce a demostrar el teorema cuando $Y$ es \emph{noetheriano}, lo que supondremos en adelante.

Podemos suponer evidentemente que $X$ e $Y$ son reducidos \sref{IV.1.5.4}.
Considerando el conjunto de las partes cerradas $Y'$ de $Y$ tales que el teorema es cierto para el subpreesquema cerrado de $Y$ cuyo espacio subyacente es $Y'$ y para $X'=f^{-1}(Y')$, podemos razonar por recurrencia noetheriana \EGAReadLink{0.2.2.2}{($\mathbf{0}_{\mathrm{I}}$,~2.2.2)} y suponer que el teorema es cierto para toda parte cerrada $Y'\neq Y$ de $Y$.
Si $X_i$ ($1\leq i\leq m$) son los subpreesquemas cerrados reducidos de $X$ cuyos espacios subyacentes son las componentes irreducibles de $X$, se tiene $\mathrm{F}_n(X)=\bigcup_i\mathrm{F}_n(X_i)$ en virtud de \sref[0]{0.14.1.5}, y podemos, por tanto, restringirnos a demostrar el teorema para cada $X_i$; en otras palabras, podemos suponer que $X$ es \emph{irreducible}.
Si $Z$ es el subpreesquema cerrado de $Y$ cuyo espacio subyacente es $\overline{f(X)}$, $f$ se factoriza como $X\xrightarrow{g}Z\xrightarrow{j}Y$, donde $j$ es la inyección canónica \sref[I]{I.5.2.2} y $g$ es de tipo finito \sref{IV.1.5.4}; basta, por tanto, demostrar el teorema para $Z$ y $g$. Por la hipótesis de recurrencia, se reduce así al caso $Z=Y$, es decir, al caso en que $Y$ es irreducible y $f$ es dominante.
Sea entonces $\eta$ el punto genérico de $Y$ y pongamos $e=\dim(f^{-1}(\eta))$; de \sref{IV.13.1.1} se sigue que, para $n\leq e$, se tiene $\mathrm{F}_n(X)=X$, y en consecuencia podemos restringirnos al caso en que $n>e$.
Pero entonces, por \sref{IV.9.5.6}, existe un entorno abierto $U$ de $\eta$ en $Y$ tal que $\mathrm{F}_n(X)\subset f^{-1}(Y-U)$; como $Y-U\neq Y$, la hipótesis de recurrencia implica que $\mathrm{F}_n(X)$ es cerrado.

II) Pasemos ahora al caso general, escribiendo todavía $S=\Spec(A)$ para el objeto meta y suponiendo $X=\Spec(B)$ afín, con $B$ una $A$-álgebra de tipo finito, luego de la forma
\oldpage[IV-3]{190}
$B=A[T_1,\ldots,T_n]/\mathfrak{J}$.
Sea $(\mathfrak{J}_\lambda)$ la familia de los ideales finitamente generados de $A[T_1,\ldots,T_n]$ contenidos en $\mathfrak{J}$, de modo que $\mathfrak{J}$ es la unión filtrante de los $\mathfrak{J}_\lambda$; si se consideran $X$ y los $X_\lambda=\Spec(A[T_1,\ldots,T_n]/\mathfrak{J}_\lambda)$ como subpreesquemas cerrados de $Z=\Spec(A[T_1,\ldots,T_n])$, se tiene, para los espacios subyacentes, $X=\bigcap_\lambda X_\lambda$.
Si $f_\lambda:X_\lambda\to S$ es el morfismo estructural, se deduce que $f^{-1}(s)=\bigcap_\lambda f_\lambda^{-1}(s)$ para todo $s\in S$; y como los conjuntos $f_\lambda^{-1}(s)$ son cerrados en el espacio noetheriano $Z_s$, existe un índice $\lambda$ (que depende de $s$) tal que $f^{-1}(s)=f_\lambda^{-1}(s)$.
Para todo $x\in X$, pongamos $d(x)=\dim_x(f^{-1}(f(x)))$ y $d_\lambda(x)=\dim_x(f_\lambda^{-1}(f_\lambda(x)))$. Lo que precede muestra que $d(x)=\inf_\lambda d_\lambda(x)$; como las funciones $d_\lambda$ son semicontinuas superiormente por la primera parte de la demostración, también lo es $d$.
\end{proof}

\begin{corollary}[13.1.4]
\label{IV.13.1.4}
Bajo las hipótesis de \sref{IV.13.1.3}, el conjunto de los $x\in X$ tales que $x$ está aislado en $f^{-1}(f(x))$ es abierto en $X$.
\end{corollary}

\begin{proof}
En efecto, es el complementario de $\mathrm{F}_1(X)$ \sref[0]{0.14.1.10}.
\end{proof}

Observemos que recuperamos así, bajo hipótesis más generales, la consecuencia \sref[III]{III.4.4.10} del «Teorema principal» de Zariski.

\begin{corollary}[13.1.5]
\label{IV.13.1.5}
Bajo las hipótesis de \sref{IV.13.1.3}, supongamos además que $f$ es un morfismo cerrado.
Entonces, para todo entero $n$, el conjunto de los $y\in Y$ tales que $\dim(f^{-1}(y))\geq n$ es cerrado; en otras palabras, la aplicación $y\mapsto\dim(f^{-1}(y))$ es semicontinua superiormente; en particular, si $y$ es una especialización de $y'$, se tiene $\dim(f^{-1}(y))\geq\dim(f^{-1}(y'))$.
\end{corollary}

\begin{proof}
En efecto, decir que $\dim(f^{-1}(y))\geq n$ significa que $y\in f(\mathrm{F}_n(X))$ \sref[0]{0.14.1.6}.
\end{proof}

\begin{corollary}[13.1.6]
\label{IV.13.1.6}
Sean $X$, $Y$ dos preesquemas irreducibles, y sea $f:X\to Y$ un morfismo dominante localmente de tipo finito.
Sea $\xi$ (resp.\ $\eta$) el punto genérico de $X$ (resp.\ de $Y$), y sea $e=\dim(f^{-1}(\eta))$.
Entonces, para todo $x\in X$, se tiene $\dim_x(f^{-1}(f(x)))\geq e$.
\end{corollary}

\begin{proof}
En efecto, el conjunto de los $x$ tales que $\dim_x(f^{-1}(f(x)))<e$ es abierto por \sref{IV.13.1.3} y, como no puede contener a $\xi$, es vacío.
\end{proof}

Observemos finalmente el siguiente resultado sencillo:

\begin{proposition}[13.1.7]
\label{IV.13.1.7}
Sea $Y$ un preesquema cuasicompacto y $f:X\to Y$ un morfismo de tipo finito.
Existe un entero $n$ tal que, para todo $y\in Y$, se tiene $\dim(f^{-1}(y))\leq n$.
\end{proposition}

\begin{proof}
Como existe un recubrimiento abierto afín finito $(V_i)$ de $Y$ tal que cada $f^{-1}(V_i)$ es una unión finita de abiertos afines, se reduce inmediatamente al caso en que $Y=\Spec(A)$ y $X=\Spec(B)$ son afines, siendo $B$ una $A$-álgebra de tipo finito.
Si $B$ admite un sistema de $n$ generadores, entonces, para todo $y\in Y$, $B\otimes_A k(y)$ es una $k(y)$-álgebra que admite $n$ generadores, luego $\dim(f^{-1}(y))\leq n$ en virtud de \sref{IV.4.1.1}.
\end{proof}

\subsection{Morfismos equidimensionales: caso de los morfismos dominantes de preesquemas irreducibles}
\label{subsection:IV.13.2}

\begin{env}[13.2.1]
\label{IV.13.2.1}
Sea $Y$ un preesquema irreducible, $X$ un preesquema irreducible, $f:X\to Y$ un morfismo \emph{dominante} localmente de tipo finito; sea $\eta$ el punto genérico de $Y$.
Sabemos por \sref{IV.13.1.6} que, para todo $x\in X$, se tiene $\dim_x(f^{-1}(f(x)))\geq\dim(f^{-1}(\eta))$.
\end{env}

\oldpage[IV-3]{191}
\begin{definition}[13.2.2]
\label{IV.13.2.2}
Bajo las hipótesis de \sref{IV.13.2.1}, se dice que $f$ es \emph{equidimensional en el punto} $x$ (o que $X$ es \emph{equidimensional sobre} $Y$ \emph{en el punto} $x$) si
\[
  \dim_x(f^{-1}(f(x)))=\dim(f^{-1}(\eta)).
\]
Se dice que $f$ es \emph{equidimensional} (o que $X$ es \emph{equidimensional sobre} $Y$) si $f$ es equidimensional en todo punto $x\in X$.
\end{definition}

Del teorema de Chevalley \sref{IV.13.1.3} se sigue que el conjunto de los $x\in X$ en los que $f$ es equidimensional es \emph{abierto} y no vacío.
Además, si $f$ es equidimensional en el punto $x$, todas las componentes irreducibles de $f^{-1}(f(x))$ que contienen a $x$ tienen la \emph{misma dimensión}, pues cada una de ellas tiene dimensión $\geq\dim(f^{-1}(\eta))$ \sref{IV.13.1.6} y $\leq\dim_x(f^{-1}(f(x)))$ en virtud de \sref[0]{0.14.1.5}.

\begin{proposition}[13.2.3]
\label{IV.13.2.3}
Sea $Y$ un preesquema irreducible localmente noetheriano, $X$ un preesquema irreducible, $f:X\to Y$ un morfismo dominante localmente de tipo finito; sea $\eta$ el punto genérico de $Y$, $x$ un punto de $X$, $y=f(x)$, y supongamos que se tiene
\[
\label{IV.13.2.3.1}
\tag{13.2.3.1}
  \dim(\sh{O}_x)=\dim(\sh{O}_y)+\dim(\sh{O}_x\otimes_{\sh{O}_y}k(y)).
\]
Entonces $f$ es equidimensional en el punto $x$.
El recíproco es cierto si se cumple la igualdad en \sref{IV.5.6.5.2}, en particular si $\sh{O}_y$ es universalmente catenario.
\end{proposition}

\begin{proof}
Esto resulta inmediatamente de \sref{IV.5.6.5.2} y de la desigualdad $\delta(x)\geq 0$ \sref{IV.13.1.1}.
\end{proof}

\begin{env}[13.2.4]
\label{IV.13.2.4}
Más generalmente, sea ahora $Y$ un preesquema localmente noetheriano, $f:X\to Y$ un morfismo de tipo finito, $x$ un punto de $X$ y $X'$ una parte cerrada irreducible de $X$ que contiene a $x$; pongamos $y=f(x)$ e $Y'=\overline{f(X')}$, y sea $\eta'$ el punto genérico de $Y'$.
Denotemos también por $X'$ e $Y'$ los subpreesquemas cerrados reducidos de $X$ e $Y$ cuyos espacios subyacentes son respectivamente $X'$ e $Y'$; la restricción $X'\to Y$ de $f$ se factoriza como $X'\xrightarrow{f'}Y'\xrightarrow{j}Y$, donde $j$ es la inyección canónica \sref[I]{I.5.2.2}, y $f'$ es de tipo finito \sref{IV.1.5.4}.
Pongamos, para abreviar,
\[
\label{IV.13.2.4.1}
\tag{13.2.4.1}
  A=\sh{O}_{Y,y},\quad B=\sh{O}_{X,x},\quad A'=\sh{O}_{Y',y},\quad B'=\sh{O}_{X',x}.
\]
La fórmula \sref{IV.5.6.5.2}, aplicada al morfismo dominante $f'$ y a los preesquemas irreducibles $X'$, $Y'$, da
\[
\label{IV.13.2.4.2}
\tag{13.2.4.2}
  \dim(B')\leq\dim(A')+\dim(B'\otimes_{A'}k(y))-(\dim_x(f'^{-1}(y))-\dim(f'^{-1}(\eta')))
\]
Por otra parte, el anillo local $A'$ (resp.\ $B'$, $B'\otimes_{A'}k(y)$) es un cociente de $A$ (resp.\ de $B$, $B\otimes_A k(y)$), luego por \sref[0]{0.16.1.2.1} se tiene
\[
\tag{13.2.4.3}
\label{IV.13.2.4.3}
  \dim(A')\leq\dim(A),\quad \dim(B')\leq\dim(B),\quad \dim(B'\otimes_{A'}k(y))\leq\dim(B\otimes_A k(y)).
\]
Se deduce, por tanto, primero de \sref{IV.13.2.4.3} y \sref[0]{0.16.3.9}
\[
\label{IV.13.2.4.4}
\tag{13.2.4.4}
  \dim(B')\leq\dim(A')+\dim(B'\otimes_{A'}k(y))\leq\dim(A)+\dim(B\otimes_A k(y)).
\]
Además, en virtud de \sref[0]{0.16.3.9}, se tienen también las desigualdades
\[
\label{IV.13.2.4.5}
\tag{13.2.4.5}
  \dim(B')\leq\dim(B)\leq\dim(A)+\dim(B\otimes_A k(y)).
\]
\end{env}

\oldpage[IV-3]{192}

Comparando estas desigualdades se obtiene:

\begin{lemma}[13.2.5]
\label{IV.13.2.5}
Con la notación de \sref{IV.13.2.4}, las condiciones siguientes son equivalentes:
\begin{enumerate}[label=\emph{\alph*)}]
  \item $\dim(B')=\dim(A)+\dim(B\otimes_A k(y))$.
  \item Se cumplen simultáneamente las relaciones siguientes:
  \begin{enumerate}[label=(\roman*)]
    \item $\dim(A')=\dim(A)$.
    \item $\dim(B'\otimes_{A'}k(y))=\dim(B\otimes_A k(y))$, en otras palabras,
    \[
      \dim_x(X'\cap f^{-1}(y))=\dim_x(f^{-1}(y)).
    \]
    \item $\dim_x(f'^{-1}(y))=\dim(f'^{-1}(\eta'))$, es decir, $f':X'\to Y'$ es equidimensional \sref{IV.13.2.2} en el punto $x$.
    \item Se tiene la igualdad
    \[
      \dim(B')=\dim(A')+\dim(B'\otimes_{A'}k(y))-(\dim_x(f'^{-1}(y))-\dim(f'^{-1}(\eta')))
    \]
  \end{enumerate}
  (relación que se cumple siempre cuando $A=\sh{O}_{Y,y}$ es un anillo universalmente catenario, en virtud de \sref{IV.5.6.5}).
  \item Se cumplen simultáneamente las relaciones siguientes:
  \begin{enumerate}[label=(\roman*)]
    \item $\dim(B')=\dim(B)$.
    \item $\dim(B)=\dim(A)+\dim(B\otimes_A k(y))$.
  \end{enumerate}
\end{enumerate}
\end{lemma}

\begin{env}[13.2.6]
\label{IV.13.2.6}
Recordemos ahora que las componentes irreducibles $X_i$ de $X$ que contienen a $x$ son en número finito y que, por \sref{IV.5.1.2.1} y \sref[0]{0.14.2.1.1},
\[
\label{IV.13.2.6.1}
\tag{13.2.6.1}
  \dim(\sh{O}_{X,x})=\sup_i\dim(\sh{O}_{X_i,x}).
\]
La equivalencia de las condiciones $b)$ y $c)$ de \sref{IV.13.2.5} implica, por tanto, teniendo en cuenta \sref[0]{0.16.3.9} y \sref[0]{0.14.2.1}:
\end{env}

\begin{proposition}[13.2.7]
\label{IV.13.2.7}
Sea $Y$ un preesquema localmente noetheriano, $f:X\to Y$ un morfismo de tipo finito, $x$ un punto de $X$, $f(x)=y$; se tiene
\[
\tag{13.2.7.1}
\label{IV.13.2.7.1}
  \dim(\sh{O}_{X,x})\leq\dim(\sh{O}_{Y,y})+\dim(\sh{O}_{X,x}\otimes_{\sh{O}_{Y,y}}k(y)).
\]
Para que los dos miembros de \sref{IV.13.2.7.1} sean iguales, es necesario y suficiente que exista una parte cerrada irreducible $X'$ de $X$ que contenga a $x$ y satisfaga simultáneamente las condiciones siguientes:
\begin{enumerate}[label=(\roman*)]
  \item Si $Y'=\overline{f(X')}$, se tiene $\dim(\sh{O}_{Y',y})=\dim(\sh{O}_{Y,y})$.
  \item $\dim_x(X'\cap f^{-1}(y))=\dim_x(f^{-1}(y))$ (en otras palabras, $X'$ contiene una de las componentes irreducibles de $f^{-1}(y)$ de dimensión maximal entre las que contienen a $x$). Esto equivale a decir que $\dim(\sh{O}_{X',x}\otimes_{\sh{O}_{Y',y}}k(y))=\dim(\sh{O}_{X,x}\otimes_{\sh{O}_{Y,y}}k(y))$.
  \item $X'$ es equidimensional sobre $Y'$ en el punto $x$; en otras palabras, se tiene
  \[
    \dim_x(X'\cap f^{-1}(y))=\dim(X'\cap f^{-1}(\eta')),
  \]
  donde $\eta'$ es el punto genérico de $Y'$ (y, en consecuencia, todas las componentes irreducibles de $X'\cap f^{-1}(y)$ que contienen a $x$ tienen la misma dimensión \sref{IV.13.2.2}).
\oldpage[IV-3]{193}
  \item Se tiene la igualdad
  \[
    \dim(\sh{O}_{X',x})=\dim(\sh{O}_{Y',y})+\dim(\sh{O}_{X',x}\otimes_{\sh{O}_{Y',y}}k(y))
  \]
  (condición siempre implicada por (iii) cuando $\sh{O}_{Y,y}$ es un anillo universalmente catenario).
\end{enumerate}
Además, $X'$ es entonces una componente irreducible de $X$, e $Y'$ una componente irreducible de $Y$.
\end{proposition}

En el enunciado, los anillos locales $\sh{O}_{X',x}$ y $\sh{O}_{Y',y}$ son relativos a los subpreesquemas cerrados reducidos de $X$, $Y$ cuyos espacios subyacentes son respectivamente $X'$ e $Y'$.

Además, esto demuestra también:

\begin{corollary}[13.2.8]
\label{IV.13.2.8}
Si los dos miembros de \sref{IV.13.2.7.1} son iguales, las partes cerradas irreducibles $X'$ de $X$ que contienen a $x$ y satisfacen las condiciones (i) a (iv) de \sref{IV.13.2.7} son exactamente aquellas para las que $\dim(\sh{O}_{X',x})=\dim(\sh{O}_{X,x})$ (lo que implica necesariamente que $X'$ es una componente irreducible de $X$).
\end{corollary}

La Proposición \sref{IV.13.2.7} implica:

\begin{corollary}[13.2.9]
\label{IV.13.2.9}
Sea $Y$ un preesquema localmente noetheriano, $f:X\to Y$ un morfismo de tipo finito, $x$ un punto de $X$, $f(x)=y$.
Supongamos que $\sh{O}_{Y,y}$ es un anillo universalmente catenario.
Entonces las condiciones siguientes son equivalentes:
\begin{enumerate}[label=\emph{\alph*)}]
  \item Los dos miembros de \sref{IV.13.2.7.1} son iguales.
  \item Existe una componente irreducible $Z$ de $f^{-1}(y)$ que contiene a $x$, de dimensión $\dim_x(f^{-1}(y))$, tal que, para todo $x'$ de un entorno de $x$ en $Z$, se tiene
  \[
  \label{IV.13.2.9.1}
\tag{13.2.9.1}
    \dim(\sh{O}_{X,x'})=\dim(\sh{O}_{Y,y})+\dim(\sh{O}_{X,x'}\otimes_{\sh{O}_{Y,y}}k(y)).
  \]
  \item Existe una componente irreducible $Z$ de $f^{-1}(y)$ que contiene a $x$, de dimensión $\dim_x(f^{-1}(y))$, tal que, para el punto genérico $z$ de $Z$, se tiene
  \[
  \label{IV.13.2.9.2}
\tag{13.2.9.2}
    \dim(\sh{O}_{X,z})=\dim(\sh{O}_{Y,y}).
  \]
\end{enumerate}
\end{corollary}

Demostremos que $a)$ implica $b)$.
Sea $\dim_x(f^{-1}(y))=n$; en virtud de \sref{IV.13.2.7}, existe una componente irreducible $X'$ de $X$ que satisface las condiciones (i) a (iv) de \sref{IV.13.2.7}; sea $Z$ una componente irreducible de dimensión $n$ de $f^{-1}(y)$, que contiene a $x$ y está contenida en $X'$.
Como $f^{-1}(y)$ es localmente noetheriano, existe un entorno abierto $U$ de $x$ en $Z$ tal que $U$ no corta ninguna componente irreducible de $f^{-1}(y)$ distinta de las que contienen a $x$, luego \sref{IV.4.1.1.3} $\dim_{x'}(f^{-1}(y))=n$ para todo $x'\in U$; entonces está claro que las condiciones (i) a (iii) de \sref{IV.13.2.7} se satisfacen cuando se sustituye $x$ por cualquier punto $x'\in U$, y lo mismo ocurre con la condición (iv), pues $\sh{O}_{Y,y}$ es universalmente catenario; de donde la conclusión por \sref{IV.13.2.7}.
La condición $b)$ implica trivialmente $c)$ por \sref{IV.5.1.2}.
Finalmente, si se satisface $c)$ y si $X''$ es una componente irreducible de $X$ que contiene a $Z$ y tal que las condiciones (i) a (iv) de \sref{IV.13.2.7} se satisfacen al sustituir $X'$ por $X''$ y $x$ por $z$, está claro que estas condiciones también se satisfacen cuando se sustituye $X'$ por $X''$ y $x$ por $x$, pues $\sh{O}_{Y,y}$ es universalmente catenario; por tanto, $c)$ implica $a)$.

\begin{proposition}[13.2.10]
\label{IV.13.2.10}
Sea $Y$ un preesquema irreducible localmente noetheriano, $\eta$ su punto genérico, $f:X\to Y$ un morfismo de tipo finito, $y$ un punto de $f(X)$.
Sean $Z_i$ las
\oldpage[IV-3]{194}
componentes irreducibles de $f^{-1}(y)$, $z_i$ el punto genérico de $Z_i$, y consideremos las condiciones siguientes:
\begin{enumerate}[label=\emph{\alph*)}]
  \item Para todo $x\in f^{-1}(y)$, se tiene la relación
  \[
  \tag{13.2.10.1}
\label{IV.13.2.10.1}
    \dim(\sh{O}_{X,x})=\dim(\sh{O}_{Y,y})+\dim(\sh{O}_{X,x}\otimes_{\sh{O}_{Y,y}}k(y)).
  \]
  \item Para todo $i$, se tiene
  \[
  \label{IV.13.2.10.2}
\tag{13.2.10.2}
    \dim(\sh{O}_{X,z_i})=\dim(\sh{O}_{Y,y}).
  \]
  \item Para todo $i$, existe una componente irreducible $X_i$ de $X$ que contiene a $Z_i$ y tal que $\dim(X_i\cap f^{-1}(\eta))=\dim(Z_i)$ (en otras palabras, el \emph{subpreesquema cerrado reducido} $X_i$ de $X$ es equidimensional sobre $Y$ en el punto $z_i$).
\end{enumerate}
Entonces $a)$ implica $b)$ y $b)$ implica $c)$; además, si $\sh{O}_{Y,y}$ es universalmente catenario, las tres condiciones $a)$, $b)$, $c)$ son equivalentes.
\end{proposition}

Como el anillo $\sh{O}_{X,z_i}\otimes_{\sh{O}_{Y,y}}k(y)$ es de dimensión $0$ \sref{IV.5.1.2}, $a)$ implica evidentemente $b)$; $b)$ implica $c)$ en virtud de \sref{IV.13.2.7}, aplicado en el punto $z_i$.
Recíprocamente, supongamos que $\sh{O}_{Y,y}$ es universalmente catenario; como la condición $c)$ implica que $f(X_i)$ es denso en $Y$, de $c)$ se sigue que las condiciones (i) a (iv) de \sref{IV.13.2.7} se satisfacen cuando se sustituye $X'$ por $X_i$ y $x$ por $z_i$, luego $c)$ implica $b)$; finalmente, $b)$ implica $a)$ por \sref{IV.13.2.9}.

\begin{corollary}[13.2.11]
\label{IV.13.2.11}
Con la notación de \sref{IV.13.2.10}, supongamos que $\sh{O}_{Y,y}$ es universalmente catenario.
Para todo $y'\in Y$, sea $E(y')$ el conjunto de las dimensiones de las componentes irreducibles de $f^{-1}(y')$ y pongamos $d(y')=\dim(f^{-1}(y'))=\sup(E(y'))$.
Entonces, si se satisfacen las condiciones equivalentes $a)$, $b)$, $c)$ de \sref{IV.13.2.10}, se tiene $E(y)\subset E(\eta)$, de donde $d(y)\leq d(\eta)$.
\end{corollary}

En efecto, con la notación de \sref{IV.13.2.10}, $X_i\cap f^{-1}(\eta)$ no es vacío y, en consecuencia, es una componente irreducible de $f^{-1}(\eta)$ \sref[0]{0.2.1.8}.

\begin{remarks}[13.2.12]
\label{IV.13.2.12}
\begin{enumerate}[label=(\roman*)]
  \item Recordemos por \sref{IV.6.1.2} que la relación \sref{IV.13.2.10.1} se satisface siempre cuando el morfismo $f$ es \emph{plano} en el punto $x$.
  \item Con la notación de \sref{IV.13.2.10}, supongamos que $\sh{O}_{Y,y}$ es universalmente catenario y, además, que el morfismo $f$ es \emph{propio}; entonces, si se satisfacen las condiciones equivalentes $a)$, $b)$, $c)$ de \sref{IV.13.2.10}, se tiene $d(y)=d(\eta)$, pues de \sref{IV.13.1.5} se sigue que $d(y)\geq d(\eta)$.
  \item El morfismo de \sref{IV.12.2.3},~(ii) es \emph{propio} y \emph{plano}, y todos los anillos locales de $Y$ son universalmente catenarios; además, las dos componentes irreducibles $X_1$, $X_2$ de $X$ son \emph{equidimensionales} sobre $Y$ en todo punto; pero $E(y)$ tiene \emph{dos} elementos para todo $y\neq y_0$, mientras que $E(y_0)$ se reduce a un \emph{solo} elemento, luego $E(y)$ \emph{no es constante} en $Y$.
\end{enumerate}
\end{remarks}

\subsection{Morfismos equidimensionales: caso general}
\label{subsection:IV.13.3}

\begin{proposition}[13.3.1]
\label{IV.13.3.1}
Sea $Y$ un preesquema, $f:X\to Y$ un morfismo localmente de tipo finito, $x$ un punto de $X$, $y=f(x)$.
Denotemos por $Y_\alpha$ las componentes irreducibles de $Y$ que contienen a $y$.
Entonces las condiciones siguientes son equivalentes:
\begin{enumerate}[label=\emph{\alph*)}]
\oldpage[IV-3]{195}
  \item Existen un entero $e$ y un entorno abierto $U$ de $x$ tales que la imagen por $f$ de toda componente irreducible de $U$ es densa en algún $Y_\alpha$, y tales que, para todo $x'\in U$, el espacio $U\cap f^{-1}(f(x'))$ es equidimensional y de dimensión $e$.
  \item[\emph{a$'$)}] Existen un entero $e$ y un entorno abierto $U$ de $x$ tales que la imagen por $f$ de toda componente irreducible de $U$ es densa en algún $Y_\alpha$, y tales que, si se denota por $y_\alpha$ el punto genérico de $Y_\alpha$, toda componente irreducible de los espacios $U\cap f^{-1}(y)$, $U\cap f^{-1}(y_\alpha)$ es de dimensión $e$.
  \item[\emph{a$''$)}] Existen un entero $e$ y un entorno abierto $U$ de $x$ tales que, para cada una de las componentes irreducibles $U_\lambda$ de $U$, $f(U_\lambda)$ es denso en algún $Y_\alpha$ y, para todo $x'\in U_\lambda$, las componentes irreducibles de $U_\lambda\cap f^{-1}(f(x'))$ son todas de dimensión $e$.
  \item Existen un entero $e$, un entorno abierto $U$ de $x$ y un $Y$-morfismo cuasifinito $g:U\to Y\otimes_{\bb{Z}}\bb{Z}[T_1,\ldots,T_e]$ (preesquema que denotaremos también $Y[T_1,\ldots,T_e]$, para abreviar) tales que la imagen por $g$ de toda componente irreducible de $U$ es densa en una componente irreducible de $Y[T_1,\ldots,T_e]$.
\end{enumerate}
\end{proposition}

\begin{proof}
Es inmediato que $a'')$ implica $a)$, pues, para todo $x\in U$, cada componente irreducible de $U\cap f^{-1}(f(x))$ es una componente irreducible de uno de los $U_\lambda\cap f^{-1}(f(x))$, de donde la conclusión por \sref[0]{0.14.1.4}.
La condición $a)$ implica trivialmente $a')$.
Demostremos que $a')$ implica $a'')$; podemos restringirnos al caso en que $f$ es de tipo finito y $X$ e $Y$ son reducidos \sref{IV.1.5.4}.
Sea $U_\lambda$ una componente irreducible de $U$, y supongamos que $f(U_\lambda)$ es denso en $Y_\alpha$; entonces la restricción de $f$ a $U_\lambda$ se factoriza como $U_\lambda\xrightarrow{f_\alpha}Y_\alpha\longrightarrow Y$, donde $f_\alpha$ es de tipo finito y dominante \sref[I]{I.5.2.2}.
Sea $x_\lambda$ el punto genérico de $U_\lambda$; en virtud de \sref[0]{0.2.1.8}, $U_\lambda\cap f_\alpha^{-1}(y_\alpha)$ es la única componente irreducible de $U\cap f^{-1}(y_\alpha)$ que contiene a $x_\lambda$ y, por hipótesis, es de dimensión $e$, igual a las dimensiones de todas las componentes irreducibles de $U_\lambda\cap f^{-1}(y)$, en virtud de la hipótesis y de \sref{IV.13.1.1}.
Pero, por el teorema de Chevalley \sref{IV.13.1.3} y \sref{IV.13.1.1}, el conjunto $V_\lambda$ de los $x'\in U_\lambda$ tales que $\dim_{x'}(f_\lambda^{-1}(f_\lambda(x')))=e$ es abierto y contiene a $x$, y basta tomar la unión de los $V_\lambda$ como conjunto abierto que satisface las condiciones de $a'')$.

Demostremos ahora que $a)$ implica $b)$; podemos restringirnos al caso en que $Y=\Spec(A)$ y $X=\Spec(B)$ son afines y $U=X$.
Demostremos primero el lema siguiente:

\begin{lemma}[13.3.1.1]
\label{IV.13.3.1.1}
Sean $Y=\Spec(A)$, $X=\Spec(B)$ dos esquemas afines, $f:X\to Y$ un morfismo de tipo finito, $y$ un punto de $f(X)$, y pongamos $e=\dim(f^{-1}(y))$.
Entonces existe un $Y$-morfismo $g:X\to Y[T_1,\ldots,T_e]$ tal que (si se pone $X_y=X\times_Y\Spec(k(y))$) el morfismo $g_y:X_y\to\Spec(k(y)[T_1,\ldots,T_e])$ es finito.
Además, para todo morfismo $g$ de este tipo, $g_y$ es necesariamente sobreyectivo, y existe un entorno abierto $U$ de $f^{-1}(y)$ en $X$ tal que $g|U$ es cuasifinito.
\end{lemma}

\begin{proof}
Sea $\mathfrak{p}=\mathfrak{j}_y$; el anillo $B\otimes_A k(\mathfrak{p})$ es una $k(\mathfrak{p})$-álgebra de tipo finito, luego, por el lema de normalización (Bourbaki, \emph{Alg.\ comm.}, cap.~V, \textsection~3, n$^{\text{o}}$~1, teo.~1), existen en $B\otimes_A k(\mathfrak{p})$ una sucesión finita $(t_i)_{1\leq i\leq r}$ de elementos algebraicamente independientes sobre $k(\mathfrak{p})$ y tales que, si se pone $C'=k(\mathfrak{p})[t_1,\ldots,t_r]$, $B\otimes_A k(\mathfrak{p})$ es una $C'$-álgebra \emph{finita}; se tiene, por tanto, $\dim(B\otimes_A k(\mathfrak{p}))=\dim(C')$ \sref[0]{0.16.1.5} y, como $\dim(C')=r$ \sref{IV.5.2.1}, se tiene $r=e$.
Como $B\otimes_A k(\mathfrak{p})=(B/\mathfrak{p}B)\otimes_{A/\mathfrak{p}}k(\mathfrak{p})$, se puede, multiplicando los $t_i$ por un elemento $\neq0$
\oldpage[IV-3]{196}
conveniente de $A/\mathfrak{p}$, suponer que cada $t_i$ es la imagen canónica en $B\otimes_A k(\mathfrak{p})$ de un elemento $s_i\in B$.
Sea entonces $u:A[T_1,\ldots,T_e]\to B$ el homomorfismo tal que $u(T_i)=s_i$ para todo $i$, y sea $g:X\to Y[T_1,\ldots,T_e]$ el morfismo correspondiente.
Está claro que, por la elección de los $s_i$, $g_y$ es un morfismo finito.
Para todo morfismo $g:X\to Y[T_1,\ldots,T_e]$ tal que $g_y$ es finito, de \sref{IV.5.4.2} y \sref{IV.4.1.2.1} se sigue que $g_y$ es necesariamente sobreyectivo.
Por otra parte, por \sref{IV.13.1.4}, el conjunto $U$ de los $x\in X$ que están aislados en su fibra $g^{-1}(g(x))$ es abierto en $X$ y contiene a $f^{-1}(y)$, y por definición la restricción $g|U$ es un morfismo cuasifinito.
\end{proof}

Establecido este lema, para demostrar que $a)$ implica $b)$ queda mostrar que, si $x_\lambda$ es un punto maximal de $U$, entonces $g(x_\lambda)=z_\lambda$ es un punto maximal de $Z=Y[T_1,\ldots,T_e]$.
Por la hipótesis $a)$, podemos (disminuyendo $U$ si es necesario) suponer que $f(x_\lambda)$ es uno de los puntos genéricos $y_\alpha$ de las componentes irreducibles de $Y$ que contienen a $y$. Si $p:Z\to Y$ es el morfismo estructural, entonces $p(z_\lambda)=y_\alpha$, y $g$ induce, por tanto, un $k(y_\alpha)$-morfismo cuasifinito
\[
h:U\cap f^{-1}(y_\alpha)\longrightarrow p^{-1}(y_\alpha).
\]
Ahora bien, $p^{-1}(y_\alpha)=\Spec(k(y_\alpha)[T_1,\ldots,T_e])$ es íntegro y de dimensión $e$. Si $z_\lambda$ no fuera un punto maximal de $Z$, tampoco sería un punto maximal de $p^{-1}(y_\alpha)$ \sref[0]{0.2.1.8}, y su adherencia $Z'_\lambda$ en $p^{-1}(y_\alpha)$ tendría, por tanto, dimensión $<e$ \sref{IV.4.1.2.1}. Pero, como $h$ es cuasifinito, de \sref{IV.4.1.2} y de la hipótesis $a)$ se sigue que $\dim(Z'_\lambda)\geq e$ (la restricción de $h$ a $\overline{\{x_\lambda\}}$ se factoriza como
\[
\overline{\{x_\lambda\}}\longrightarrow Z'_\lambda\longrightarrow p^{-1}(y_\alpha)
\]
por \sref[I]{I.5.2.2}). Esto es una contradicción y demuestra que $a)$ implica $b)$.

Demostremos finalmente que $b)$ implica $a)$.
Obsérvese que el morfismo estructural $p:Z=Y[T_1,\ldots,T_e]\to Y$ es fielmente plano; por tanto, según \sref{IV.2.3.4}, las imágenes por $p$ de los puntos maximales de $Z$ son los puntos maximales de $Y$. Esto demuestra ya que los $f(x_\lambda)$ son puntos genéricos de los $Y_\alpha$.
Además, si $f(x_\lambda)=y_\alpha$, el $k(y_\alpha)$-morfismo $h:U\cap f^{-1}(y_\alpha)\to p^{-1}(y_\alpha)$ deducido de $g$ es dominante y cuasifinito por hipótesis; por tanto, según \sref{IV.4.1.2}, se tiene $\dim(U_\lambda\cap f^{-1}(y_\alpha))=e$ (siendo $U_\lambda$ la componente irreducible de $U$ de punto genérico $x_\lambda$).
Análogamente, para todo $x'\in U_\lambda$, el morfismo $U_\lambda\cap f^{-1}(f(x'))\to p^{-1}(f(x'))$ deducido de $g$ es un $k(f(x'))$-morfismo cuasifinito, luego por \sref{IV.4.1.2} se tiene $\dim(f^{-1}(f(x'))\cap U_\lambda)\leq e$; pero, por \sref{IV.13.1.6}, todas las componentes irreducibles de $U_\lambda\cap f^{-1}(f(x'))$ son de dimensión $\geq e$; se ve, por tanto, que estas componentes son exactamente de dimensión $e$, luego $b)$ implica $a)$.
\end{proof}

\begin{definition}[13.3.2]
\label{IV.13.3.2}
Sea $Y$ un preesquema, $f:X\to Y$ un morfismo localmente de tipo finito, $x$ un punto de $X$.
Se dice que $f$ es \emph{equidimensional en el punto} $x$ (o que $X$ es \emph{equidimensional sobre} $Y$ \emph{en el punto} $x$) si se satisfacen las condiciones equivalentes de \sref{IV.13.3.1}.
Se dice que $f$ es \emph{equidimensional} (o que $X$ es \emph{equidimensional sobre} $Y$) si $f$ es equidimensional en todo punto $x\in X$.
\end{definition}

Está claro que, para que $f$ sea equidimensional en un punto $x$, es necesario y suficiente que $f_{\mathrm{red}}$ lo sea.
Además, las condiciones de \sref{IV.13.3.1} muestran que el conjunto de los puntos en los que $f$ es equidimensional es \emph{abierto} en $X$.

Observemos que, cuando $X$ e $Y$ son \emph{irreducibles}, decir que $f$ es equidimen-
\oldpage[IV-3]{197}
sional en el punto $x$ significa, en virtud de \sref{IV.13.3.1},~$a'')$, que $f$ es dominante y que $\dim_x(f^{-1}(f(x)))=\dim(f^{-1}(\eta))$ (donde $\eta$ es el punto genérico de $Y$); la definición \sref{IV.13.3.2} coincide así, en este caso, con la definición \sref{IV.13.2.2}.

\begin{proposition}[13.3.3]
\label{IV.13.3.3}
Sea $Y$ un preesquema localmente noetheriano, $f:X\to Y$ un morfismo localmente de tipo finito, $x$ un punto de $X$ y $X_j$ las componentes irreducibles de $X$ (en número finito) que contienen a $x$.
Para que $f$ sea equidimensional en el punto $x$, es necesario y suficiente que, para todo $j$, $Y_j=\overline{f(X_j)}$ sea una componente irreducible de $Y$ y que, denotando por $X_j$ e $Y_j$ los subpreesquemas cerrados reducidos de $X$ e $Y$ de espacios subyacentes $X_j$ e $Y_j$, por $f_j:X_j\to Y_j$ el morfismo deducido de $f$ \sref[I]{I.5.2.2}, y por $y_j$ el punto genérico de $Y_j$, $f_j$ sea equidimensional en el punto $x$ y todos los números $\dim(f_j^{-1}(y_j))$ sean iguales.
\end{proposition}

\begin{proof}
Esto resulta inmediatamente de \sref{IV.13.3.1},~$a')$).
\end{proof}

\begin{corollary}[13.3.4]
\label{IV.13.3.4}
Con la notación de \sref{IV.13.3.3}, pongamos $y=f(x)$.
Si $\sh{O}_{Y,y}$ es un anillo universalmente catenario y $f$ es equidimensional en el punto $x$, se tiene
\[
\tag{13.3.4.1}
\label{IV.13.3.4.1}
  \dim(\sh{O}_{X_j,x})=\dim(\sh{O}_{Y_j,y})+e-\deg.\mathrm{tr}_{k(y)}k(x)
\]
donde $e$ es el valor común de los números $\dim(f_j^{-1}(y_j))$.
\end{corollary}

\begin{proof}
Como cada uno de los $\sh{O}_{Y_j,y}$, cociente de $\sh{O}_{Y,y}$, es un anillo universalmente catenario \sref{IV.5.6.1}, la igualdad es consecuencia de \sref{IV.5.6.5}.
\end{proof}

\begin{corollary}[13.3.5]
\label{IV.13.3.5}
Con la notación de \sref{IV.13.3.4}, supongamos que $f$ es equidimensional en el punto $x$ y que $\sh{O}_{Y,y}$ es un anillo universalmente catenario.
Si el anillo $\sh{O}_{Y,y}$ es equidimensional, también lo es $\sh{O}_{X,x}$; el recíproco es cierto si la imagen por $f$ de la unión de los $X_j$ es densa en un entorno de $y$.
\end{corollary}

\begin{proof}
En efecto, decir que $\sh{O}_{Y,y}$ es equidimensional significa que los números $\dim(\sh{O}_{Y'_i,y})$ son iguales para todas las componentes irreducibles $Y'_i$ de $Y$ que contienen a $y$, como resulta de \sref{IV.5.1.1.5}, aplicado al esquema local $\Spec(\sh{O}_{Y,y})$; como se tiene la relación \sref{IV.13.3.4.1}, del mismo razonamiento se sigue que $\sh{O}_{X,x}$ es entonces equidimensional.
Recíprocamente, si $\sh{O}_{X,x}$ es equidimensional, todos los números $\dim(\sh{O}_{Y_j,y})$ son iguales por \sref{IV.13.3.4.1}; se deduce que $\sh{O}_{Y,y}$ es equidimensional si los $Y_j$ son todas las componentes irreducibles de $Y$ que contienen a $y$, lo cual resulta de la hipótesis suplementaria.
\end{proof}

\begin{proposition}[13.3.6]
\label{IV.13.3.6}
Sea $Y$ un preesquema localmente noetheriano, $f:X\to Y$ un morfismo localmente de tipo finito, $x$ un punto de $X$, $y=f(x)$.
Supongamos que el anillo $\sh{O}_y$ es equidimensional. Entonces, si $\sh{O}_x$ es equidimensional y se tiene la igualdad
\[
\tag{13.3.6.1}
\label{IV.13.3.6.1}
  \dim(\sh{O}_x)=\dim(\sh{O}_y)+\dim(\sh{O}_x\otimes_{\sh{O}_y}k(y))
\]
(cf.\ \sref{IV.13.2.7.1}), $f$ es equidimensional en el punto $x$, y el recíproco es cierto si $\sh{O}_y$ es un anillo universalmente catenario.
\end{proposition}

\begin{proof}
Conservando la notación de \sref{IV.13.3.3}, de \sref{IV.13.2.8} y de la hipótesis de que $\sh{O}_x$ es equidimensional se sigue que cada $Y_j$ es una componente irreducible de $Y$ y que cada $f_j$ es equidimensional en el punto $x$; además, por \sref{IV.13.2.8}, se tiene (teniendo en cuenta \sref{IV.5.6.5})
\[
  \dim(\sh{O}_{X_j,x})=\dim(\sh{O}_{Y_j,y})+\dim_x(X_j\cap f^{-1}(y))-\deg.\mathrm{tr}_{k(y)}k(x).
\]

\oldpage[IV-3]{198}

Ahora bien, como $\sh{O}_y$ se supone equidimensional, esta igualdad puede escribirse
\[
\tag{13.3.6.2}
\label{IV.13.3.6.2}
  \dim_x(X_j\cap f^{-1}(y))=\dim(\sh{O}_{X,x})-\dim(\sh{O}_{Y,y})+\deg.\mathrm{tr}_{k(y)}k(x).
\]
El primer miembro de \sref{IV.13.3.6.2} es, por tanto, independiente de $j$; pero, como $f_j$ es equidimensional en el punto $x$, se tiene $\dim_x(X_j\cap f^{-1}(y))=\dim(f_j^{-1}(y_j))$, luego el criterio de \sref{IV.13.3.3} muestra que $f$ es equidimensional en el punto $x$.

Recíprocamente, supongamos que $\sh{O}_{Y,y}$ es un anillo universalmente catenario y que $f$ es equidimensional en el punto $x$; entonces, por \sref{IV.13.3.5}, $\sh{O}_{X,x}$ es equidimensional, y de \sref{IV.13.3.4} se sigue que se tiene la relación
\[
  \dim(\sh{O}_{X,x})=\dim(\sh{O}_{Y,y})+e-\deg.\mathrm{tr}_{k(y)}k(x)
\]
donde $e$ es el valor común de los números $\dim(f_j^{-1}(y_j))=\dim_x(X_j\cap f^{-1}(y))$; pero, por definición, $e$ es también igual a $\dim_x(f^{-1}(y))$, de donde la relación \sref{IV.13.3.6.1}, teniendo en cuenta \sref{IV.5.6.5.2}.
\end{proof}

\begin{proposition}[13.3.7]
\label{IV.13.3.7}
Sea $Y$ un preesquema, $f:X\to Y$ un morfismo localmente de tipo finito y equidimensional, y $Z$ una parte cerrada de $X$.
Entonces la función
\[
\label{IV.13.3.7.1}
\tag{13.3.7.1}
  x\mapsto\operatorname{codim}_x(Z\cap f^{-1}(f(x)),f^{-1}(f(x)))
\]
es semicontinua inferiormente en $X$.
\end{proposition}

\begin{proof}
Como todas las componentes irreducibles de $f^{-1}(f(x))$ que contienen a $x$ tienen la misma dimensión por hipótesis \sref{IV.13.3.1}, se tiene, en virtud de \sref{IV.5.2.1},
\[
  \operatorname{codim}_x(Z\cap f^{-1}(f(x)),f^{-1}(f(x)))=\dim_x(f^{-1}(f(x)))-\dim_x(Z\cap f^{-1}(f(x))).
\]
Pero, por hipótesis, el primer término del segundo miembro es una función continua de $x$ \sref{IV.13.3.1}, y el segundo es una función semicontinua superiormente de $x$ por \sref{IV.13.1.3}\footnote{La fuente impresa cita aquí (13.3.3); la semicontinuidad superior invocada es el teorema de Chevalley (13.1.3).}; de donde la conclusión.
\end{proof}

\begin{proposition}[13.3.8]
\label{IV.13.3.8}
Sea $Y$ un preesquema, $f:X\to Y$ un morfismo localmente de tipo finito, $x$ un punto de $X$.
Sea $Y'$ un preesquema, $g:Y'\to Y$ un morfismo plano, $X'=X\times_Y Y'$, $f'=f_{(Y')}:X'\to Y'$; si $f$ es equidimensional en el punto $x$, entonces $f'$ es equidimensional en todo punto $x'\in X'$ situado sobre $x$.
\end{proposition}

\begin{proof}
Como la cuestión es local sobre $X$ e $Y$, podemos restringirnos al caso en que toda componente irreducible de $X$ (resp.\ de $Y$) contiene a $x$ (resp.\ a $y$); como la imagen por $f$ de toda componente irreducible de $X$ es densa en una componente irreducible de $Y$ \sref{IV.13.3.1}, se sabe por \sref{IV.2.3.5} que la imagen por $f'$ de toda componente irreducible de $X'$ es densa en una componente irreducible de $Y'$.
Por otra parte, por transitividad de las fibras \sref[I]{I.3.6.4}, de \sref{IV.4.2.8} y \sref{IV.13.3.1} se sigue que, para todo $z'\in X'$, el conjunto de las dimensiones de las componentes irreducibles de ${f'}^{-1}(f'(z'))$ es el mismo que el conjunto de las dimensiones de las componentes irreducibles de $f^{-1}(f(z))$, donde $z$ es la proyección de $z'$ sobre $X$; de donde la conclusión en virtud de \sref{IV.13.3.1},~$a)$).
\end{proof}

\begin{remarks}[13.3.9]
\label{IV.13.3.9}
La propiedad de equidimensionalidad de un morfismo $f:X\to Y$ \emph{no} es estable por un cambio de base arbitrario $g:Y'\to Y$, ni siquiera cuando $g$ es la inyección canónica de una componente irreducible $Y'$ de $Y$ en $Y$ (siendo $Y'$
\oldpage[IV-3]{199}
considerado como subpreesquema cerrado reducido de $Y$).
Por ejemplo, sea $k$ un cuerpo, sea $A_0=k[S,T]$ el anillo de polinomios en dos indeterminadas, y pongamos $\mathfrak{a}=\mathfrak{p}_1\mathfrak{p}_2$, donde $\mathfrak{p}_1$ y $\mathfrak{p}_2$ son los ideales primos $A_0S$ y $A_0T$ de $A_0$; sea $A=A_0/\mathfrak{a}$ e $Y=\Spec(A)$, que tiene dos componentes irreducibles $Y_1=\Spec(A/\mathfrak{p}_1)$ e $Y_2=\Spec(A/\mathfrak{p}_2)$; tomemos $X=Y_1$, con $f:X\to Y$ la inyección canónica, que es evidentemente un morfismo equidimensional.
Por otra parte, tomemos $Y'=Y_2$, con $g:Y'\to Y$ la inyección canónica; entonces $X'=X\times_Y Y'=\Spec\bigl((A/\mathfrak{p}_1)\otimes_A(A/\mathfrak{p}_2)\bigr)=\Spec\bigl(A/(\mathfrak{p}_1+\mathfrak{p}_2)\bigr)$, y $\mathfrak{p}_1+\mathfrak{p}_2$ es un ideal maximal de $A$, de modo que $X'$ se reduce a un punto; el morfismo $f':X'\to Y'$ \emph{no es dominante}, pues la imagen por $f'$ del único punto de $X'$ es un punto cerrado de $Y'$; en consecuencia, $f'$ no es equidimensional.

También se puede dar un contraejemplo en el que $X$ e $Y$ son \emph{íntegros}, $f$ es \emph{finito} y \emph{birracional} (y \emph{a fortiori} equidimensional por \sref{IV.13.3.1},~$b)$), y $g:Y'\to Y$ es \emph{finito} y \emph{dominante}.
Sean $A$ y $\overline{A}$ los anillos locales definidos en \sref{IV.11.7.5}, y tomemos $Y=\Spec(A)$ y $X=\Spec(\overline{A})$; por otra parte, con la notación de \sref{IV.11.7.5}, tomemos $Y'=\Spec(\overline{B})$; entonces $X'=\Spec(\overline{A}\otimes_A\overline{B})=\Spec(\overline{B}\otimes_B\overline{B})$; pero se verifica inmediatamente que $\overline{B}\otimes_B\overline{B}$ es el producto directo de los anillos $B/\mathfrak{p}'$, $B/\mathfrak{p}''$ y de dos anillos isomorfos a $B/\mathfrak{n}$, cuyos espectros se reducen, por tanto, a un punto; como las proyecciones de estos puntos son puntos cerrados de $Y'$, se ve aquí de nuevo que $f'$ no lleva una componente irreducible de $X'$ sobre una parte densa de una componente irreducible de $Y'$, y por tanto $f'$ no es equidimensional.

En este ejemplo, el anillo $A$ no es geométricamente unirramificado; veremos en el parágrafo siguiente {\hyperlink{ega4.c04.cor.14-4-6}{(14.4.6)}} que tales fenómenos no pueden producirse cuando los puntos de $Y$ son \emph{geométricamente unirramificados}.
La falta de estabilidad de la noción de morfismo equidimensional restringe mucho su interés, en favor de la noción de morfismo universalmente abierto, que se estudiará en detalle en el parágrafo siguiente.
\end{remarks}
